\section{Заключение}

В рамках практики решены две взаимосвязанные задачи: разработка системы
анализа производительности для эмулятора CUDA PTX и верификация
корректности эмулятора на реальных CUDA-ядрах.

\textbf{Система профилирования.} Реализован профилировщик с четырьмя
метриками (\texttt{branch\_efficiency}, \texttt{bank\_conflicts},
\texttt{global\_coalescing}, \texttt{write\_ops}), интегрированный в
эмулятор через кодогенерацию по JSON-манифесту. Архитектурное решение~---
хранение метаданных метрик в \texttt{profiling\_metrics.json} с автогенерацией
C++-кода скриптом \texttt{autogen.py}~--- позволяет добавлять новые метрики
без модификации горячего пути исполнения. Буферизация записей
\texttt{ProfilingRecord} на уровне варпа с единственным сбросом в общий
профилировщик по завершении блока минимизирует накладные расходы на синхронизацию.
Инструмент генерации HTML-отчёта (\texttt{renderer.py}) формирует
самодостаточный файл с пятью секциями, обеспечивая переход от сводных
метрик к аннотированному исходному коду в рамках одного документа.

\textbf{Верификация корректности.} Все семь интеграционных тестов пройдены
при максимальном отклонении от CPU-эталона не более $10^{-3}$.
Охваченные вычислительные паттерны: поэлементные операции (\texttt{vadd}),
трансцендентные функции (\texttt{gelu}), многофазные редукции (\texttt{softmax}),
тайловое матричное умножение (\texttt{mmul}), Flash Attention-подобное ядро,
FFT и итерационный метод сопряжённых градиентов. Верификация системы
профилирования подтвердила нулевые отклонения по всем шести аналитическим
тестам метрик.

\textbf{Производительность.} Замедление относительно V100-PCIE составляет
1\,900--2\,300$\times$ для нетривиальных нагрузок~--- неизбежная стоимость
программной эмуляции без JIT-компиляции. Накладные расходы профилировщика
при включённом сборе данных составляют 88--130\% для лёгких ядер и
500--800\% для вычислительно плотных (тайловое матричное умножение),
что приемлемо для задач анализа и отладки.

Среди текущих ограничений~--- неполнота покрытия PTX~ISA: ряд инструкций
не реализован вовсе (тензорные операции \texttt{mma} и \texttt{ldmatrix}),
а у реализованных поддержаны не все варианты модификаторов. Эмулятор
прошёл проверку на учебных примерах и лабораторных работах, однако
остаётся первой версией инструмента и требует существенного расширения
покрытия для применения на произвольных промышленных кодах.

\newpage
